\documentclass[10pt, aspectratio=169]{beamer}

\usepackage[spanish]{babel}
\usepackage[utf8]{inputenc}
\usepackage[pscoord]{eso-pic}
\usepackage{hyperref}
\usepackage{eqnarray,amsmath}
\usepackage[absolute,overlay]{textpos}
\hypersetup{colorlinks=true,urlcolor=blue}

\usepackage{tikz}
\usetikzlibrary{automata, positioning, arrows.meta, petri}

\usepackage{listings}
\usepackage{xcolor}

\usepackage[style=apa]{biblatex}
\bibliography{biblio.bib}

\title{}
\author{}
\date{}

\newtheorem{teorema}{Teorema.}
\setbeamersize{text margin left=10pt,text margin right=10pt}
\beamertemplatenavigationsymbolsempty

\hypersetup{
	colorlinks=true,
	pdfborderstyle={/S/U/W 1}
}

\setbeamertemplate{background}{\includegraphics[width=\paperwidth,height=\paperheight,keepaspectratio]{portada.pdf}}

\newcommand{\azulULA}{\color[rgb]{0.30078125,0.30078125,0.999999}}
\definecolor{azulULA}{RGB}{77, 77, 255}

\setbeamercolor{title}{fg=azulULA}
\setbeamercolor{frametitle}{fg=azulULA}
\setbeamercolor{block title}{bg=azulULA,fg=white}
\setbeamercolor{block body}{bg=azulULA!10!bg}

\renewcommand{\qedsymbol}{\textcolor{black}{\openbox}}

\newcommand{\placetextbox}[3]{
	\setbox0=\hbox{#3}
	\AddToShipoutPictureFG*{
		\put(\LenToUnit{#1\paperwidth},\LenToUnit{#2\paperheight}){\vtop{{\null}\makebox[0pt][c]{#3}}}
	}
}
\newcommand{\placetextboxleft}[3]{
	\setbox0=\hbox{#3}
	\AddToShipoutPictureFG*{
		\put(\LenToUnit{#1\paperwidth},\LenToUnit{#2\paperheight}){\vtop{{\null}\makebox[0pt][l]{#3}}}
	}
}

\begin{document}
	% Diapositiva de portada
	\begin{frame}
		\placetextbox{0.47}{0.40}{\large\azulULA\textbf{\begin{tabular}{c} Teoría de la Computación y Lenguajes Formales: \\
					Sesión 12-B. El Motor en Acción: Planificación, Supervisión y Diagnóstico \end{tabular}}}
	\end{frame}
	
	% Cambio de formato para el cuerpo
	\setbeamertemplate{background}{
		\includegraphics[width=\paperwidth,height=\paperheight,keepaspectratio]{fondo1280x720.pdf}}
	
	\begin{frame}{Alcance de la sesión 12-B}
		\tableofcontents
	\end{frame}

% =======================================================
\section{El Motor de Eventos Discretos como Gemelo Digital}
% =======================================================

\begin{frame}{}
	\centering
	\Huge{El Motor de Eventos Discretos como Gemelo Digital: \\ De la Red Verificada al Run-Time}
\end{frame}

\begin{frame}{El Ciclo de Vida Completo: ¿Dónde se ubica la Sesión 12-B?}
	\begin{block}{De la Especificación a la Operación Real}
		Tras superar las compuertas de sintaxis (PDA - Clase 10), completitud ontológica y verificación formal de \emph{Soundness} (Clase 12-A), entramos al **Paso 8**:
	\end{block}

	\begin{center}
		\begin{tikzpicture}[x=3.4cm, y=-1.6cm, >=Stealth, thick]
			% Fila 1: Diseño, Sintaxis y Semántica previa
			\node[draw=gray!60, rectangle, rounded corners, fill=gray!10, text width=2.6cm, align=center] (p1) at (0, 0)
				{\tiny 1. Entender dominio};
			\node[draw=gray!60, rectangle, rounded corners, fill=gray!10, text width=2.6cm, align=center] (p2) at (1, 0)
				{\tiny 2. Modelos RdP};
			\node[draw=gray!60, rectangle, rounded corners, fill=gray!10, text width=2.6cm, align=center] (p3) at (2, 0)
				{\tiny 3. Codificar RdP+Onto};
			\node[draw=gray!60, rectangle, rounded corners, fill=gray!10, text width=2.6cm, align=center] (p4) at (3, 0)
				{\tiny 4. Validar PDA (Clase 10)};

			% Fila 2: Verificación previa y Ejecución actual
			\node[draw=gray!60, rectangle, rounded corners, fill=gray!10, text width=2.6cm, align=center] (p5) at (3, 1)
				{\tiny 5. Completitud Ontología};
			\node[draw=gray!60, rectangle, rounded corners, fill=gray!10, text width=2.6cm, align=center] (p6) at (2, 1)
				{\tiny 6. Verificar Soundness (12-A)};
			\node[draw=gray!60, rectangle, rounded corners, fill=gray!10, text width=2.6cm, align=center] (p7) at (1, 1)
				{\tiny 7. Desplegar al Motor};
			\node[draw=azulULA, rectangle, rounded corners, fill=azulULA!25, text width=2.6cm, align=center, line width=1.5pt] (p8) at (0, 1)
				{\textbf{8. Usar Sistema}\\ \scriptsize \textbf{Planificar, Supervisar,}\\ \textbf{Diagnosticar (12-B)}};

			\draw[->, thick, gray!60] (p1) -- (p2);
			\draw[->, thick, gray!60] (p2) -- (p3);
			\draw[->, thick, gray!60] (p3) -- (p4);
			\draw[->, thick, gray!60] (p4) -- (p5);
			\draw[->, thick, gray!60] (p5) -- (p6);
			\draw[->, thick, gray!60] (p6) -- (p7);
			\draw[->, thick, azulULA, line width=1.2pt] (p7) -- (p8);
		\end{tikzpicture}
	\end{center}
\end{frame}

\begin{frame}{La Arquitectura del Motor (Fénix): Del Modelo a la Ejecución}
	\begin{block}{El Motor de Redes de Petri como Núcleo de Ejecución}
		Una vez verificado el modelo formal (Clase 12-A), el motor carga la matriz de incidencia $C$, los lugares de recursos y los retardos nominales $\tau$:
	\end{block}

	\begin{columns}
		\column{.50\textwidth}
		\begin{block}{Entradas del Mundo Físico / Negocio}
			\begin{itemize}
				\item Señales de sensores industriales (PLC, RFID, IoT).
				\item Triggers de usuarios (aprobación en ERP, clic en botón de UI).
				\item Reloj del sistema (tiempos transcurridos).
			\end{itemize}
		\end{block}

		\column{.48\textwidth}
		\begin{alertblock}{Funciones del Motor Fénix}
			\begin{itemize}
				\item Mantiene el estado en memoria ($M_k$).
				\item Dispara transiciones válidas según reglas formales.
				\item Conecta tres fases operacionales: \textbf{Planificar}, \textbf{Supervisar} y \textbf{Diagnosticar}.
			\end{itemize}
		\end{alertblock}
	\end{columns}
\end{frame}

\begin{frame}{El Marcado $M_k$ como Estado del Gemelo Digital}
	\begin{block}{La Foto Viva de la Planta y del Negocio}
		En cualquier instante de tiempo $t$, el vector de marcado $M(t)$ describe con rigor matemático:
		\begin{itemize}
			\item \textbf{Dónde está cada producto o trámite:} Fichas en lugares de proceso (\texttt{Pieza\_En\_Mecanizado}, \texttt{Orden\_En\_Revision}).
			\item \textbf{Qué recursos están libres u ocupados:} Fichas en lugares de recursos (\texttt{Robot\_1\_Libre}, \texttt{Analista\_Compras\_Disponible}).
			\item \textbf{Qué pasos están habilitados inmediatamente:} El conjunto de transiciones activas:
			\[ T_{\text{hab}}(M_k) = \{ t \in T \mid M_k \ge {}^\bullet t \} \]
		\end{itemize}
	\end{block}

	\begin{alertblock}{Sincronización Continua}
		El gemelo digital no es una animación 3D pasiva; es un **modelo ejecutable** que evoluciona en paralelo exacto con la realidad mediante la ecuación de estado:
		\[ M_{k+1} = M_k + C \cdot u_k \]
	\end{alertblock}
\end{frame}

% =======================================================
\section{Fase 1: Planificación y Programación (Pre-ejecución)}
% =======================================================

\begin{frame}{}
	\centering
	\Huge{Fase 1: Planificación y Programación: \\ Simulación, Lead Time y Ruta Crítica}
\end{frame}

\begin{frame}{Redes Temporizadas para el Cálculo del \emph{Lead Time}}
	\begin{block}{De la Causalidad al Tiempo Métrico}
		Asociando duraciones nominales $\tau(t)$ a las transiciones (T-timed) o tiempos de permanencia $\tau(p)$ a los lugares (P-timed), el motor simula el flujo antes de iniciarlo:
	\end{block}

	\begin{columns}
		\column{.52\textwidth}
		\begin{itemize}
			\item \textbf{Tiempo de Ciclo Teórico (\emph{Makespan}):} Tiempo mínimo necesario para que una ficha viaje desde la fuente $i$ hasta el sumidero $o$.
			\item \textbf{Compromiso Real con el Cliente:} Permite cotizar fechas de entrega basadas en la capacidad disponible de los lugares de recursos, no en suposiciones.
			\item \textbf{Evaluación de Escenarios (\emph{What-if}):} ¿Qué pasa si ingresan 10 pedidos simultáneos? ¿Qué pasa si una máquina opera a media velocidad?
		\end{itemize}

		\column{.45\textwidth}
		\begin{center}
			\begin{tikzpicture}[node distance=1.4cm, >=Stealth, thick]
				\node[place, tokens=1] (i) [label={above:\scriptsize $i$}] {};
				\node[transition] (t1) [right of=i, label={above:\scriptsize $t_1 (\tau=10)$}] {} edge [pre] (i);
				\node[place] (p1) [right of=t1, label={above:\scriptsize $p_1$}] {} edge [pre] (t1);
				\node[transition] (t2) [right of=p1, label={above:\scriptsize $t_2 (\tau=25)$}] {} edge [pre] (p1);
				\node[place] (o) [right of=t2, label={above:\scriptsize $o$}] {} edge [pre] (t2);
			\end{tikzpicture}
		\end{center}
		\centering{\footnotesize $\text{Lead Time Nominal} = 10 + 25 = 35\,\text{min}$.}
	\end{columns}
\end{frame}

\begin{frame}{Ruta Crítica y Cuellos de Botella Estructurales}
	\begin{columns}
		\column{.48\textwidth}
		\begin{center}
			\begin{tikzpicture}[node distance=1.25cm, >=Stealth, thick]
				\node[place, tokens=1] (i) [label={left:\tiny $i$}] {};
				\node[transition] (split) [right of=i, label={above:\tiny split}] {} edge [pre] (i);
				
				\node[place] (pA) [above right of=split, yshift=0.1cm, label={above:\tiny Rápida}] {} edge [pre] (split);
				\node[transition] (tA) [right of=pA, label={above:\tiny $t_A (\tau=5\text{s})$}] {} edge [pre] (pA);
				\node[place] (pA2) [right of=tA, label={above:\tiny Listo A}] {} edge [pre] (tA);
				
				\node[place] (pB) [below right of=split, yshift=-0.1cm, label={below:\tiny Lenta}] {} edge [pre] (split);
				\node[transition] (tB) [right of=pB, label={below:\tiny $t_B (\tau=40\text{s})$}] {} edge [pre] (pB);
				\node[place] (pB2) [right of=tB, label={below:\tiny Listo B}] {} edge [pre] (tB);
				
				\node[transition] (join) [below right of=pA2, yshift=0.35cm, label={above:\tiny join}] {}
					edge [pre] (pA2) edge [pre] (pB2);
				\node[place] (o) [right of=join, label={right:\tiny $o$}] {} edge [pre] (join);
			\end{tikzpicture}
		\end{center}

		\column{.48\textwidth}
		\begin{alertblock}{Ruta Crítica y Espera Pasiva}
			\small
			\begin{itemize}
				\item La \textbf{Ruta Crítica} está dictada por la rama lenta $t_B$ (40 s).
				\item La ficha de la Rama A esperará $35\,\text{s}$ ociosa en el lugar de sincronización $pA_2$.
				\item \textbf{Detección Previa de WIP:} El motor identifica antes de ejecutar dónde se acumulará inventario en proceso (\emph{Work in Process}).
			\end{itemize}
		\end{alertblock}
	\end{columns}
\end{frame}

% =======================================================
\section{Fase 2: Supervisión y Control en Tiempo Real}
% =======================================================

\begin{frame}{}
	\centering
	\Huge{Fase 2: Supervisión y Control en Tiempo Real: \\ Triggers, Watchdogs y Alarmas Tempranas}
\end{frame}

\begin{frame}{Seguimiento en Vivo: Triggers y Desambiguación por Estado}
	\begin{block}{El Despachador de Eventos de Fénix}
		Durante la operación, el motor recibe un flujo constante de eventos externos:
		\[ \text{Evento entrante: } \langle \text{trigger} = \texttt{'FIN\_TAREA'}, \, \text{timestamp} = 10\text{:}14\text{:}02 \rangle \]
	\end{block}

	\begin{columns}
		\column{.50\textwidth}
		\begin{block}{Mecanismo de Despacho}
			\begin{enumerate}
				\item Fénix busca entre las transiciones activas:
				\[ T_{\text{cand}} = \{ t \in T_{\text{hab}}(M_k) \mid \ell(t) = \text{trigger} \} \]
				\item Si $|T_{\text{cand}}| = 1$: Dispara la transición $t$ y actualiza el marcado:
				\[ M_{k+1} = M_k + C \cdot u \]
			\end{enumerate}
		\end{block}

		\column{.48\textwidth}
		\begin{alertblock}{Diagnóstico en Vivo de Anomalías}
			\begin{itemize}
				\item Si $|T_{\text{cand}}| = 0$: \textbf{¡Alerta de incoherencia físico-digital!}
				\item Llegó una señal de un sensor en un estado donde era imposible que ocurriera.
				\item Permite detectar piezas saltadas, sensores averiados o pulsaciones dobles de operadores.
			\end{itemize}
		\end{alertblock}
	\end{columns}
\end{frame}

\begin{frame}{Temporizadores Activos (\emph{Watchdogs}) y Detección Temprana}
	\begin{block}{Vigilancia Continua de Tiempos Límite}
		En cada estado o actividad se establece una cota máxima: $T_{\max}(p) = \tau_{\text{nominal}}(p) + \Delta_{\text{tol}}$.
	\end{block}

	\begin{columns}
		\column{.48\textwidth}
		\begin{center}
			\begin{tikzpicture}[node distance=1.3cm, >=Stealth, thick]
				\node[draw=azulULA, rectangle, rounded corners, fill=azulULA!10, text width=3.4cm, align=center] (in) 
					{\textbf{Entrada al Estado}\\ \scriptsize Ficha entra en $p$ / Cronómetro ON};
					
				\node[draw=green!60!black, rectangle, rounded corners, fill=green!10, text width=3.4cm, align=center, below of=in, yshift=-0.1cm] (ok) 
					{\textbf{Disparo Nominal}\\ \scriptsize Trigger antes de $T_{\max}$ $\implies$ OK};
					
				\node[draw=red!80, rectangle, rounded corners, fill=red!10, text width=3.4cm, align=center, below of=ok, yshift=-0.1cm] (al) 
					{\textbf{Expiración Watchdog}\\ \scriptsize Supera $T_{\max}$ $\implies$ \alert{Alarma Demora}};

				\draw[->, thick] (in) -- (ok);
				\draw[->, thick] (ok) -- (al);
			\end{tikzpicture}
		\end{center}

		\column{.48\textwidth}
		\begin{alertblock}{Ventaja Operativa Inmediata}
			\begin{itemize}
				\item El supervisor no espera al final de la jornada para enterarse de la demora.
				\item **El gemelo digital alerta mientras el retraso está sucediendo**, habilitando acciones correctivas tempranas antes de que afecte a los clientes.
			\end{itemize}
		\end{alertblock}
	\end{columns}
\end{frame}

% =======================================================
\section{Fase 3: Diagnóstico Causal de Demoras (Post-ejecución)}
% =======================================================

\begin{frame}{}
	\centering
	\Huge{Fase 3: Diagnóstico Causal de Demoras: \\ ¿Por qué un estado demoró más de lo esperado?}
\end{frame}

\begin{frame}{El Problema del Diagnóstico: ¿Quién es el culpable?}
	\begin{block}{La Pregunta Clave de Rendimiento}
		Al auditar el log de eventos, una orden tardó $90\,\text{minutos}$ en vez de los $30\,\text{minutos}$ planificados.
		\textbf{¿Por qué demoró el estado?}
	\end{block}

	\begin{columns}
		\column{.32\textwidth}
		\begin{alertblock}{Hipótesis 1: Recursos}
			¿El recurso (máquina, operario, vehículo) estaba ocupado, averiado o ausente?
		\end{alertblock}

		\column{.32\textwidth}
		\begin{alertblock}{Hipótesis 2: Insumos}
			¿Los insumos eran de mala calidad, fuera de medida o defectuosos, obligando a retrabajos?
		\end{alertblock}

		\column{.32\textwidth}
		\begin{alertblock}{Hipótesis 3: Operación}
			¿La persona o máquina tardó más en realizar la actividad que su valor nominal?
		\end{alertblock}
	\end{columns}

	\vspace{0.3cm}
	\begin{exampleblock}{El Poder del Modelo Formal}
		La Red de Petri desacopla explícitamente el **Recurso**, el **Insumo/Producto** y la **Transición**, permitiendo aislar la causa raíz con certeza matemática.
	\end{exampleblock}
\end{frame}

\begin{frame}{Causa 1: Fricción por Recursos (Cola vs. Servicio)}
	\begin{block}{Descomposición Formal del Tiempo de Permanencia}
		Para una pieza o trámite en un estado, el tiempo total observado $T_{\text{total}}$ se compone de:
		\[ T_{\text{total}} = T_{\text{espera}} (\text{Cola por Recursos}) + T_{\text{servicio}} (\text{Transformación Activa}) \]
	\end{block}

	\begin{columns}
		\column{.52\textwidth}
		\begin{center}
			\begin{tikzpicture}[node distance=1.8cm, >=Stealth, thick]
				\node[place, tokens=1] (cola) [label={above:\scriptsize $p_{\text{cola}}$ (Pieza lista)}] {};
				\node[place, tokens=0] (rec) [below of=cola, yshift=0.3cm, label={below:\scriptsize $p_{\text{rec}}$ (\alert{Recurso Ocupado})}] {};
				\node[transition] (t) [right of=cola, yshift=-0.75cm, label={above:\scriptsize $t_{\text{iniciar}}$}] {}
					edge [pre] (cola) edge [pre] (rec);
				\node[place] (proc) [right of=t, label={above:\scriptsize $p_{\text{proceso}}$}] {} edge [pre] (t);
			\end{tikzpicture}
		\end{center}

		\column{.45\textwidth}
		\begin{alertblock}{Diagnóstico por Marcado de Recurso}
			\begin{itemize}
				\item Si $M(p_{\text{cola}}) \ge 1$ pero $M(p_{\text{rec}}) = 0$:
				\item La tarea **no puede arrancar** porque el recurso está capturado en otro proceso o en mantenimiento.
				\item \textbf{Causa Raíz:} Cuello de botella por escasez de recursos, no problema de la pieza ni del operador.
			\end{itemize}
		\end{alertblock}
	\end{columns}
\end{frame}

\begin{frame}{Causa 2: Calidad de Insumos y Retrabajos (Tokens CPN)}
	\begin{block}{El Rastro de la Calidad en Redes Coloreadas}
		Cuando el insumo presenta variaciones dimensionales, dureza anormal o defectos:
	\end{block}

	\begin{columns}
		\column{.50\textwidth}
		\begin{block}{Efecto A: Rutas Alternas de Retrabajo}
			\begin{itemize}
				\item El control de calidad dispara una transición de desvío:
				\[ \text{token} = \{ \text{lote: } \text{'L-40'}, \, \text{calidad: } \text{'Defectuoso'} \} \]
				\item La ficha recorre un lazo de reproceso:
				\[ p_{\text{ensamble}} \to t_{\text{rechazo}} \to p_{\text{reparar}} \to p_{\text{ensamble}} \]
				\item \textbf{Diagnóstico:} El estado tardó más porque **se ejecutó 2 o más veces**. Causa: proveedor o lote previo.
			\end{itemize}
		\end{block}

		\column{.48\textwidth}
		\begin{block}{Efecto B: Sobreesfuerzo de Maquinado}
			\begin{itemize}
				\item En materiales más duros, los sensores de par/potencia del husillo reducen el avance.
				\item La duración de la transición se dilata:
				\[ \tau_{\text{real}}(t) \gg \tau_{\text{nominal}}(t) \]
				\item \textbf{Diagnóstico:} Inconsistencia en la materia prima suministrada.
			\end{itemize}
		\end{block}
	\end{columns}
\end{frame}

\begin{frame}{Matriz de Diagnóstico Causal del Motor}
	\begin{block}{Algoritmo de Identificación de Causa Raíz en Fénix}
		\small
		\begin{tabular}{|p{2.8cm}|p{3.0cm}|p{2.9cm}|p{3.3cm}|}
			\hline
			\textbf{Síntoma} & \textbf{Estado del Marcado} & \textbf{Atributos del Token} & \textbf{Causa Raíz Determinada} \\ \hline\hline
			Retardo antes de iniciar la actividad & $M(p_{\text{pieza}}) \ge 1$ \newline $M(p_{\text{recurso}}) = 0$ & Irrelevante & \textbf{Falta de Recurso} (contención o máquina parada). \\ \hline
			La actividad se ejecutó varias veces & Se disparó transición de desvío & \texttt{estado = 'No\_OK'} \newline \texttt{reintentos > 0} & \textbf{Calidad de Insumo} (retrabajo forzado). \\ \hline
			La actividad inició a tiempo pero tardó más & $M(p_{\text{recurso}}) \ge 1$ durante todo el lapso & Parámetros nominales de pieza & \textbf{Variabilidad Operativa} (desgaste, lentitud humana). \\ \hline
		\end{tabular}
	\end{block}
\end{frame}

% =======================================================
\section{Caso de Estudio Integrador}
% =======================================================

\begin{frame}{}
	\centering
	\Huge{Caso de Estudio Integrador: \\ Celda de Mecanizado, Recurso Compartido e Inspección}
\end{frame}

\begin{frame}{Caso Práctico: Diagnóstico en una Celda de Producción}
	\begin{block}{Estructura del Proceso}
		Una pieza llega a $p_{\text{in}}$. Requiere un Torno CNC compartido ($p_{\text{torno}}$). Al terminar, pasa por Inspección: si está OK va a $p_{\text{fin}}$, si falla va a $p_{\text{retrabajo}}$:
	\end{block}

	\begin{center}
		\begin{tikzpicture}[node distance=1.9cm, >=Stealth, thick]
			\node[place, tokens=1] (piezas) [label={above:\scriptsize $p_{\text{in}}$ (Pieza)}] {};
			\node[place, tokens=1] (torno) [below of=piezas, yshift=-0.1cm, label={below:\scriptsize $p_{\text{torno}}$ (CNC Libre)}] {};
			
			\node[transition] (t_mec) [right of=piezas, yshift=-1.0cm, label={above:\scriptsize $t_{\text{torneado}} (\tau=15\text{m})$}] {}
				edge [pre] (piezas) edge [pre] (torno);
				
			\node[place] (p_insp) [right of=t_mec, label={above:\scriptsize $p_{\text{insp}}$}] {} edge [pre] (t_mec);
			\path (t_mec) edge [post, bend left=30] (torno);
			
			\node[transition] (t_ok) [above right of=p_insp, yshift=-0.2cm, label={above:\scriptsize $t_{\text{ok}}$}] {} edge [pre] (p_insp);
			\node[place] (p_fin) [right of=t_ok, label={above:\scriptsize $p_{\text{fin}}$ (Listo)}] {} edge [pre] (t_ok);
			
			\node[transition] (t_defect) [below right of=p_insp, yshift=0.2cm, label={below:\scriptsize $t_{\text{defect}}$}] {} edge [pre] (p_insp);
			\node[place] (p_ret) [below of=t_mec, yshift=-0.4cm, label={below:\scriptsize $p_{\text{retrabajo}}$}] {};
			\path (t_defect) edge [post] (p_ret);
			\path (p_ret) edge [post, bend left=25] (t_mec);
		\end{tikzpicture}
	\end{center}
\end{frame}

\begin{frame}{Auditoría de Dos Casos Reales en el Motor Fénix}
	\begin{columns}
		\column{.48\textwidth}
		\begin{alertblock}{Caso 1: Pieza TK-801 (Demora: 65 min)}
			\begin{itemize}
				\item \textbf{Log del Gemelo Digital:}
				\begin{itemize}
					\item Llegada a $p_{\text{in}}$: 08:00
					\item $t_{\text{torneado}}$ dispara: 08:45 (\alert{45 min en espera})
					\item Fin de torneado: 09:00 (15 min nominal)
					\item $t_{\text{ok}}$ dispara: 09:05 $\implies$ Salida.
				\end{itemize}
				\item \textbf{Diagnóstico del Motor:}
				$M(p_{\text{torno}}) = 0$ entre 08:00 y 08:45 por lote previo urgente.
				\item \textbf{Causa Raíz:} **Falta de Recurso (Contención de Máquina)**.
			\end{itemize}
		\end{alertblock}

		\column{.48\textwidth}
		\begin{exampleblock}{Caso 2: Pieza TK-802 (Demora: 50 min)}
			\begin{itemize}
				\item \textbf{Log del Gemelo Digital:}
				\begin{itemize}
					\item Llegada a $p_{\text{in}}$: 09:10
					\item Torneado 1 arranca: 09:10 (0 min espera)
					\item Inspección a 09:25: \alert{Dispara $t_{\text{defect}}$}
					\item Torneado 2 arranca: 09:30 (retrabajo)
					\item Fin y salida conforme: 09:50.
				\end{itemize}
				\item \textbf{Diagnóstico del Motor:}
				La pieza recorrió el lazo de reproceso.
				\item \textbf{Causa Raíz:} **Calidad de Insumo Defectuosa**.
			\end{itemize}
		\end{exampleblock}
	\end{columns}
\end{frame}

% =======================================================
\section{Conclusión del Ciclo de Vida Completo}
% =======================================================

\begin{frame}{Conclusión de las Semanas 11 y 12}
	\begin{block}{La Gran Síntesis: De la Teoría de la Computación a la Industria 4.0}
		\small
		\begin{itemize}
			\item \textbf{Semana 11-A:} Redes de Petri como lenguajes formales y concurrencia ($a^n b^n c^n$, ecuación de estado).
			\item \textbf{Semana 11-B:} Subclases, temporización, tríptica holónica y desambiguación por marcado en Fénix.
			\item \textbf{Semana 12-A:} Verificación formal previa (\emph{Soundness}, acotamiento, invariantes, sifones y trampas).
			\item \textbf{Semana 12-B:} El motor en acción:
			\begin{enumerate}
				\small
				\item **Planificar:** Simulación de tiempos de ciclo y ruta crítica antes de producir.
				\item **Supervisar:** Gemelo digital en tiempo real alimentado por triggers de sensores.
				\item **Diagnosticar:** Descomposición formal de demoras: ¿Recursos disputados o Calidad de insumos?
			\end{enumerate}
		\end{itemize}
	\end{block}

	\begin{exampleblock}{Mensaje Central para el Ingeniero}
		\small
		Las Redes de Petri no son solo un dibujo: **son una especificación matemática ejecutable que permite diseñar sin fallos, supervisar con precisión y auditar causas con rigor científico**.
	\end{exampleblock}
\end{frame}

% Diapositiva final
\setbeamertemplate{headline}{}
\setbeamertemplate{background}
{\includegraphics[width=\paperwidth,height=\paperheight,keepaspectratio]{salida.pdf}}

\begin{frame}
	\placetextbox{0.6}{0.87}{\Huge\textbf{¡Muchas gracias!}}
\end{frame}

\setbeamertemplate{headline}{}
\setbeamertemplate{background}{}

\printbibliography
\end{document}
